% GENERATED FILE. Edit canonical lessons, then run npm run generate:books.
% canonical-source-hash: fnv1a64:cbe11ae100c7e226
% canonical-lessons: AR-C256-rasaba, AR-C256-ijtaza, AR-C256-sharaha, AR-C256-tarjama, AR-C256-aada

\chapter{To fail, To pass, To explain, To translate, To give back}
\label{ch:ar-to-fail-to-pass-to-explain-to-translate-to-give-back}
\hlchaptermodality{256}

\begin{chapteropening}
\textbf{By the end of this chapter:} \emph{I can say to fail, to pass, to explain, to translate and to give back.}

Complete the last lesson of chapter 256: I can say to fail, to pass, to explain, to translate and to give back.
\end{chapteropening}

\section[\texorpdfstring{\ar{رسب} (rasaba)}{رسب (rasaba)}]{\ar{رسب} (rasaba) — to fail (an exam)}

\label{lesson:AR-C256-rasaba}



\begin{quote}
Before the new one: say the Arabic for to get better, then the Arabic for to examine.
\end{quote}

\subsection*{You'll want to know: \ar{رسب}}
\textbf{\ar{رسب}} — \emph{rasaba} — \textquotedblleft{}to fail (an exam)\textquotedblright{}.

\begin{cousinweb}
\textbf{\ar{رسب}} (\emph{rasaba}) is built on the native Arabic root \emph{r-s-b}, \textquotedblleft{}to sink\textquotedblright{} (Lane). \emph{Rasaba} is \textquotedblleft{}to sink to the bottom\textquotedblright{}, hence \textquotedblleft{}to fail\textquotedblright{}.
\end{cousinweb}

\begin{grammarlens}[title={What you've built}]
One more word for this chapter.
\end{grammarlens}

\subsection*{Your turn}
\begin{itemize}
\raggedright
  \item \emph{Say it:} \emph{rasaba}
  \item \emph{Say it:} \emph{rasaba}, once more
  \item \emph{Say it:} say \emph{rasaba}
  \item \emph{Recall:} say the Arabic for to get better, then the Arabic for to examine, then say \emph{rasaba} again
\end{itemize}

\begin{tcolorbox}[breakable,colback=teal!4,colframe=teal!35!black,title={Before you move on}]
What does \ar{رسب} mean? (\textquotedblleft{}To fail (an exam)\textquotedblright{}.) Say it once more.
\end{tcolorbox}
\section[\texorpdfstring{\ar{اجتاز} (ijtāza)}{اجتاز (ijtāza)}]{\ar{اجتاز} (ijtāza) — to pass (an exam)}

\label{lesson:AR-C256-ijtaza}



\begin{quote}
Before the new one: say the Arabic for to examine, then the Arabic for to fail.
\end{quote}

\subsection*{You'll want to know: \ar{اجتاز}}
\textbf{\ar{اجتاز}} — \emph{ijtāza} — \textquotedblleft{}to pass (an exam)\textquotedblright{}.

\begin{cousinweb}
\textbf{\ar{اجتاز}} (\emph{ijtāza}) is built on the native Arabic root \emph{j-w-z}, \textquotedblleft{}to pass, to be allowed\textquotedblright{} (Lane). \emph{Ijtāza} is Form VIII, \textquotedblleft{}to pass through\textquotedblright{}.
\end{cousinweb}

\begin{grammarlens}[title={What you've built}]
One more word for this chapter.
\end{grammarlens}

\subsection*{Your turn}
\begin{itemize}
\raggedright
  \item \emph{Say it:} \emph{ijtāza}
  \item \emph{Say it:} \emph{ijtāza}, once more
  \item \emph{Say it:} say \emph{ijtāza}
  \item \emph{Recall:} say the Arabic for to examine, then the Arabic for to fail, then say \emph{ijtāza} again
\end{itemize}

\begin{tcolorbox}[breakable,colback=teal!4,colframe=teal!35!black,title={Before you move on}]
What does \ar{اجتاز} mean? (\textquotedblleft{}To pass (an exam)\textquotedblright{}.) Say it once more.
\end{tcolorbox}
\section[\texorpdfstring{\ar{شرح} (sharaḥa)}{شرح (sharaḥa)}]{\ar{شرح} (sharaḥa) — to explain}

\label{lesson:AR-C256-sharaha}



\begin{quote}
Before the new one: say the Arabic for to fail, then the Arabic for to pass.
\end{quote}

\subsection*{You'll want to know: \ar{شرح}}
\textbf{\ar{شرح}} — \emph{sharaḥa} — \textquotedblleft{}to explain\textquotedblright{}.

\begin{cousinweb}
\textbf{\ar{شرح}} (\emph{sharaḥa}) is built on the native Arabic root \emph{sh-r-ḥ}, \textquotedblleft{}to open, to lay open\textquotedblright{} (Lane). \emph{Sharaḥa} is \textquotedblleft{}to lay open\textquotedblright{}, to explain.
\end{cousinweb}

\begin{grammarlens}[title={What you've built}]
One more word for this chapter.
\end{grammarlens}

\subsection*{Your turn}
\begin{itemize}
\raggedright
  \item \emph{Say it:} \emph{sharaḥa}
  \item \emph{Say it:} \emph{sharaḥa}, once more
  \item \emph{Say it:} say \emph{sharaḥa}
  \item \emph{Recall:} say the Arabic for to fail, then the Arabic for to pass, then say \emph{sharaḥa} again
\end{itemize}

\begin{tcolorbox}[breakable,colback=teal!4,colframe=teal!35!black,title={Before you move on}]
What does \ar{شرح} mean? (\textquotedblleft{}To explain\textquotedblright{}.) Say it once more.
\end{tcolorbox}
\section[\texorpdfstring{\ar{ترجم} (tarjama)}{ترجم (tarjama)}]{\ar{ترجم} (tarjama) — to translate}

\label{lesson:AR-C256-tarjama}



\begin{quote}
Before the new one: say the Arabic for to pass, then the Arabic for to explain.
\end{quote}

\subsection*{You'll want to know: \ar{ترجم}}
\textbf{\ar{ترجم}} — \emph{tarjama} — \textquotedblleft{}to translate\textquotedblright{}.

\begin{cousinweb}
\textbf{\ar{ترجم}} (\emph{tarjama}) is the four-letter verb that Lane gives with \emph{tarjumān}, \textquotedblleft{}interpreter\textquotedblright{}. \emph{Tarjumān} is usually traced to Aramaic \emph{targmānā}; English \emph{dragoman} comes from it.
\end{cousinweb}

\begin{grammarlens}[title={What you've built}]
One more word for this chapter.
\end{grammarlens}

\subsection*{Your turn}
\begin{itemize}
\raggedright
  \item \emph{Say it:} \emph{tarjama}
  \item \emph{Say it:} \emph{tarjama}, once more
  \item \emph{Say it:} say \emph{tarjama}
  \item \emph{Recall:} say the Arabic for to pass, then the Arabic for to explain, then say \emph{tarjama} again
\end{itemize}

\begin{tcolorbox}[breakable,colback=teal!4,colframe=teal!35!black,title={Before you move on}]
What does \ar{ترجم} mean? (\textquotedblleft{}To translate\textquotedblright{}.) Say it once more.
\end{tcolorbox}
\section[\texorpdfstring{\ar{أعاد} (aʿāda)}{أعاد (aʿāda)}]{\ar{أعاد} (aʿāda) — to give back, to return (something)}

\label{lesson:AR-C256-aada}



\begin{quote}
Before the new one: say the Arabic for to explain, then the Arabic for to translate.
\end{quote}

\subsection*{You'll want to know: \ar{أعاد}}
\textbf{\ar{أعاد}} — \emph{aʿāda} — \textquotedblleft{}to give back, to return (something)\textquotedblright{}.

\begin{cousinweb}
\textbf{\ar{أعاد}} (\emph{aʿāda}) is built on the native Arabic root \emph{ʿ-w-d}, \textquotedblleft{}to return\textquotedblright{} (Lane). \emph{Aʿāda} is Form IV, \textquotedblleft{}to make return\textquotedblright{}.
\end{cousinweb}

\begin{grammarlens}[title={What you've built}]
One more word for this chapter.
\end{grammarlens}

\subsection*{Your turn}
\begin{itemize}
\raggedright
  \item \emph{Say it:} \emph{aʿāda}
  \item \emph{Say it:} \emph{aʿāda}, once more
  \item \emph{Say it:} say \emph{aʿāda}
  \item \emph{Recall:} say the Arabic for to explain, then the Arabic for to translate, then say \emph{aʿāda} again
\end{itemize}

\begin{tcolorbox}[breakable,colback=teal!4,colframe=teal!35!black,title={Before you move on}]
What does \ar{أعاد} mean? (\textquotedblleft{}To give back, to return (something)\textquotedblright{}.) Say it once more.
\end{tcolorbox}
